\documentclass[11pt,a4paper]{report}
\usepackage[utf8]{inputenc}
\usepackage[T1]{fontenc}
\usepackage{lmodern}
\usepackage{geometry}
\usepackage{hyperref}
\usepackage{amsmath,amssymb}
\usepackage{booktabs}
\usepackage{longtable}
\usepackage{array}
\usepackage{xcolor}
\usepackage{listings}

\geometry{margin=2.5cm}
\hypersetup{colorlinks=true,linkcolor=blue,urlcolor=blue,citecolor=blue}

\lstdefinestyle{code}{
  basicstyle=\ttfamily\small,
  frame=single,
  columns=fullflexible,
  keepspaces=true,
  showstringspaces=false,
  breaklines=true,
  keywordstyle=\color{blue!70!black},
  commentstyle=\color{green!40!black},
}

\title{Sanaverkko \\ \large Installation Mode --- Implementation Documentation}
\author{Project Documentation}
\date{March 2026}

\begin{document}
\maketitle
\tableofcontents

% ---------------------------------------------------------------------------
\chapter{Overview}

\textit{Installation mode} is a set of features added to Sanaverkko that
allow the system to be operated in a live-performance or gallery-installation
context.  In this mode:

\begin{itemize}
  \item A full-screen pygame canvas displays the word network for an audience.
  \item A strip of virtual potentiometers at the bottom of the canvas lets a
        performer or the audience adjust algorithm parameters in real time by
        dragging the knobs.
  \item A seed-sentence bar lets a user enter a starting sentence that seeds
        the word network, bypassing the automatic generation for one cycle.
  \item All of these controls are saved to and restored from JSON preset files
        together with the rest of the algorithm parameters.
\end{itemize}

The implementation lives entirely in \texttt{sanaVerkkoCore.py} and was
merged into \texttt{main} from the branch
\texttt{controls-for-the-audience} on 26 March 2026.

% ---------------------------------------------------------------------------
\chapter{Audience Potentiometer Strip}

\section{Concept}

Any of the 128 slots defined in \texttt{AUDIENCE\_PARAMS} can be
\emph{exposed} to the audience.  An exposed slot appears as a knob in the
potentiometer strip rendered at the bottom of the pygame window.  Turning the
knob writes the mapped parameter value into \texttt{self.params} on every
frame, overriding whatever the wxPython control panel shows.

\section{Parameter table}

\texttt{AUDIENCE\_PARAMS} is a list of 128 tuples:

\begin{lstlisting}[style=code]
(param_key, lo, hi, label)
\end{lstlisting}

The first 44 entries cover the primary algorithm parameters (logic, audio,
rhythm, compressor/ADSR, visual).  The remaining 84 slots wrap around the
same list cyclically, giving extra reach without hard-coding redundant entries.

\section{Audience Parameters window}

The wxPython control panel has an \textbf{Audience Parameters\ldots} button
that opens a separate scrollable window.  Each row in that window corresponds
to one AUDIENCE\_PARAMS slot and contains:

\begin{itemize}
  \item a label showing the parameter name,
  \item an \textbf{Expose} checkbox to show the knob in the pygame strip,
  \item an \textbf{Invert} checkbox to reverse the knob direction (so fully
        clockwise maps to the minimum value rather than the maximum).
\end{itemize}

\section{Potentiometer strip rendering}

The strip occupies a fixed-height band (\texttt{\_POT\_H = 80\,px}) at the
very bottom of the pygame canvas.  Each exposed knob is drawn as:

\begin{itemize}
  \item a circular arc knob with a 270\textdegree{} sweep (225\textdegree{}
        at min, $-45$\textdegree{} at max),
  \item a fill bar beneath the knob showing the current fraction of the range,
  \item a 7-character abbreviated label below the bar,
  \item a small numeric value at the very bottom of the strip.
\end{itemize}

Scroll arrows appear at the left and right edges when there are more exposed
knobs than fit horizontally (slot width $= 44\,\text{px}$).  The user can
drag the strip horizontally or use the mouse wheel to scroll.

\subsection{Tooltip}

When the mouse hovers over a knob, the full parameter name is displayed in a
small semi-transparent tooltip panel positioned just above the seed bar.  The
tooltip surface is cached and only re-rendered when focus moves to a different
knob.

\section{Knob interaction}

\begin{longtable}{>{\ttfamily}p{0.22\textwidth} p{0.72\textwidth}}
\toprule
\textbf{Action} & \textbf{Effect} \\
\midrule
\endfirsthead
\toprule
\textbf{Action} & \textbf{Effect} \\
\midrule
\endhead
Click + drag up/down & Increase/decrease the raw pot value (0--127).
  Sensitivity: $\approx 0.8$ units per pixel. \\
Mouse wheel over strip & Scroll visible knobs left/right. \\
Click + drag on empty strip area & Pan the knob row horizontally. \\
\bottomrule
\end{longtable}

% ---------------------------------------------------------------------------
\chapter{Seed Sentence Bar}

\section{Concept}

The seed sentence bar sits immediately above the potentiometer strip
(\texttt{\_SEED\_H = 72\,px}).  It allows a user to type one or more words
that immediately replace the current network nodes, seeding the generative
process with a specific starting point.

\section{Workflow}

\begin{enumerate}
  \item The user clicks anywhere on the seed bar.
  \item A borderless \texttt{wx.Frame} overlay appears, positioned pixel-
        precisely over the seed bar using the SDL2 window position
        (\texttt{pygame.\_sdl2.video.Window.position}).
  \item The overlay contains a \texttt{wx.TextCtrl} pre-filled with the
        current draft (if any).  The text field has immediate keyboard focus.
  \item The user types words separated by spaces.  Only characters present in
        \texttt{gematria\_table} (a–z, å, ä, ö) are accepted; all others are
        silently filtered on commit.
  \item Pressing \textbf{Enter} commits the sentence; pressing \textbf{Esc}
        cancels.  The overlay is then destroyed (not merely hidden — see
        §\ref{sec:overlay-destroy}).
\end{enumerate}

\section{Why a wx overlay instead of pygame keyboard input}
\label{sec:focus}

On macOS, when a wxPython application owns the Cocoa event loop, clicking
the SDL/pygame window does not reliably steal keyboard focus from the wx
control panel.  Characters typed by the user go to the invisible wx text
field rather than to pygame's event queue, so no \texttt{KEYDOWN} events
ever arrive.

Embedding a \texttt{wx.TextCtrl} in a borderless overlay frame that floats
over the pygame window solves the problem: wx always wins the focus race, so
we lean into it.  The TextCtrl uses the same input machinery as the parameter
panel fields that were already known to work.

\section{Overlay destruction}
\label{sec:overlay-destroy}

The overlay is \texttt{Destroy()}ed (not \texttt{Hide()}d) when the user
commits or cancels.  A merely hidden \texttt{STAY\_ON\_TOP} \texttt{wx.Frame}
remains in the macOS Cocoa window responder chain and double-routes keystrokes
to whichever text field is active in the main parameter window, causing each
typed character to appear twice (e.g.\ typing \texttt{64} produces
\texttt{6644}).  Destruction removes the window from the Cocoa layer
completely.

\section{Committed sentence display}

After a successful commit the sentence is shown as a line of text rendered
directly onto the pygame canvas, left-aligned just above the seed bar.  This
text persists until the next commit and serves as a visible reminder to the
performer of the current seed.

\section{Chip display}

The committed words are also shown as coloured \emph{chips} in the seed bar
itself.  Each chip's hue is derived from the word's gematria value (the sum
of the Hebrew-style letter values of its characters), scaled against the
maximum gematria in the current draft.  Chip text surfaces are cached
per word to avoid redundant font rendering.  The cache is cleared whenever
the draft changes.

\section{Simulation pause during input}

While the seed overlay is open (\texttt{\_seed\_active = True}), the main
simulation loop resets \texttt{last\_process\_time} every frame.  This means
the logic worker does not fire immediately when the overlay closes, preventing
a burst computation that would stall the event loop and make the seed bar
appear unresponsive to subsequent activations.

\texttt{\_commit\_seed\_sentence()} also resets \texttt{last\_process\_time}
directly, giving one full \texttt{process\_interval} of lag-free time after
each commit regardless of how the overlay was closed.

% ---------------------------------------------------------------------------
\chapter{Preset Integration}

\section{Saving}

When the user saves a preset via \textbf{Save Preset}, three additional
arrays are appended to the JSON file:

\begin{lstlisting}[style=code]
preset_data["audience_values"]   = list(self.pot_values)    # 128 ints 0-127
preset_data["audience_exposed"]  = list(self.pot_exposed)   # 128 bools
preset_data["audience_inverted"] = list(self.pot_inverted)  # 128 bools
\end{lstlisting}

\section{Loading}

When a preset is loaded, the three arrays are read back with element-wise
coercion (missing entries default to safe values).  If the Audience
Parameters window is currently open, its rows are refreshed immediately so
the checkboxes reflect the loaded state.

% ---------------------------------------------------------------------------
\chapter{Network View Adjustment}

Before this feature branch the word-network circle was centred on the full
window height, causing nodes to overlap with the pot strip and seed bar.

The \texttt{makeWordCircle()} method now computes a reduced centre:

\begin{lstlisting}[style=code]
pot_h    = self._POT_H if self._pot_exposed_indices() else 0
reserved = pot_h + self._SEED_H
centre_y = (self.size[1] - reserved) / 2
\end{lstlisting}

This shifts the network upward so it occupies only the area above the
audience controls, keeping nodes clearly visible at all zoom levels.

% ---------------------------------------------------------------------------
\chapter{Technical Reference}

\section{Key constants}

\begin{longtable}{>{\ttfamily}p{0.38\textwidth} p{0.56\textwidth}}
\toprule
\textbf{Name} & \textbf{Value / meaning} \\
\midrule
\endfirsthead
\toprule
\textbf{Name} & \textbf{Value / meaning} \\
\midrule
\endhead
\_POT\_H & 80 px — height of the potentiometer strip \\
\_POT\_W & 44 px — width per knob slot \\
\_POT\_R & 14 px — knob circle radius \\
\_SEED\_H & 72 px — height of the seed sentence bar \\
\_SEED\_FONT\_SIZE & 19 pt — chip label font size \\
AUDIENCE\_PARAMS & Module-level list of 128 \texttt{(key, lo, hi, label)} tuples \\
\bottomrule
\end{longtable}

\section{Key methods}

\begin{longtable}{>{\ttfamily}p{0.38\textwidth} p{0.56\textwidth}}
\toprule
\textbf{Method} & \textbf{Responsibility} \\
\midrule
\endfirsthead
\toprule
\textbf{Method} & \textbf{Responsibility} \\
\midrule
\endhead
\_pot\_exposed\_indices() & Return sorted list of pot indices where
  \texttt{pot\_exposed[i]} is True. \\
\_pot\_value\_to\_param(idx, raw) & Map raw 0--127 value to the physical
  parameter range, applying inversion. \\
\_apply\_pot\_values\_to\_params() & Called every frame; writes all exposed
  pot values into \texttt{self.params}. \\
\_draw\_pot\_strip(screen) & Render all visible knobs, scroll arrows, and
  the hover tooltip. \\
\_handle\_pot\_events(events) & Process mouse drag/scroll on knobs; update
  hover index. \\
\_draw\_seed\_editor(screen) & Render the seed bar background, chips, and
  (when overlay active) a dim overlay tint. \\
\_handle\_seed\_events(events) & Detect click on seed bar; call
  \texttt{wx.CallAfter(\_show\_seed\_overlay)}. \\
\_show\_seed\_overlay() & Create and display the borderless wx input frame. \\
\_hide\_seed\_overlay(cancel, text) & Destroy the overlay; on commit, parse
  text and call \texttt{\_commit\_seed\_sentence()}. \\
\_commit\_seed\_sentence() & Replace network nodes with seed words; reset
  audio signature; delay next logic tick. \\
\_seed\_bar\_rect() & Return \texttt{(x, y, w, h)} of the seed bar in pygame
  window coordinates. \\
\_seed\_overlay\_screen\_rect() & Translate seed bar rect to absolute screen
  coordinates using SDL2 window position. \\
\_seed\_recompute\_gem\_max() & Recompute the maximum gematria value across
  the current draft for chip colour scaling. \\
OnOpenAudienceParams(event) & Open the Audience Parameters scrollable
  wxPython window. \\
OnSavePreset(event) & Save all params plus audience arrays to JSON. \\
OnLoadPreset(event) & Restore params and audience arrays from JSON. \\
\bottomrule
\end{longtable}

\section{Data flow}

\begin{enumerate}
  \item \textbf{Every simulation frame:}
        \begin{enumerate}
          \item \texttt{\_handle\_pot\_events} processes mouse events and
                updates \texttt{pot\_values[i]}.
          \item \texttt{\_handle\_seed\_events} checks for a click on the
                seed bar and posts \texttt{\_show\_seed\_overlay} if needed.
          \item \texttt{\_draw\_seed\_editor} and \texttt{\_draw\_pot\_strip}
                render the bottom UI bands.
        \end{enumerate}
  \item \textbf{On each rate-limited logic tick} (governed by
        \texttt{process\_interval}):
        \begin{enumerate}
          \item \texttt{\_apply\_pot\_values\_to\_params} writes knob values
                into \texttt{self.params}.
          \item The logic worker reads \texttt{self.params} for all simulation
                decisions.
        \end{enumerate}
  \item \textbf{On seed commit:}
        \begin{enumerate}
          \item \texttt{\_commit\_seed\_sentence} replaces \texttt{self.words}
                and \texttt{self.referenceWords}.
          \item Audio signature is invalidated so \texttt{updateAudio()}
                re-synthesises on the next tick without interrupting the
                currently playing audio.
        \end{enumerate}
\end{enumerate}

\end{document}
